% ==============================================================================
% PDF-2 — Structural Diagnosability & Phase Logic
% Section 06a: Falsifiability by Example — Counterexample Templates (Epistemic)
%
% Purpose:
%   Provide explicit counterexample templates (logical refuters) showing how core
%   proposition types could be falsified under admissible assumptions.
%
% Constraints:
%   - No new primitives
%   - No prescriptive language
%   - No algorithms / operational procedures
%   - No empirical content
%   - Epistemic falsification only (existence-of-construction statements)
% ==============================================================================

\subsection{Falsifiability by example: counterexample templates (epistemic)}
\label{subsec:falsifiability-by-example}

This subsection gives \emph{counterexample templates} in a strict epistemic
sense. Each template specifies:
(i) admissible assumptions,
(ii) a construction class,
(iii) which proposition type would be refuted \emph{if} such a construction
exists.
No operational testing or search procedure is implied.

\paragraph{CE-1 (Non-separability under the adopted observational interface).}
\begin{enumerate}
  \item \textbf{Assumptions.}
  Fix the observational regime $(E,\Pi)$ and the induced notion of admissible
  observation traces $D \subseteq \Sigma_{obs}(K_{EA},E)$ as used in
  \Cref{sec:pdf2_structural_diagnosability}. Assume the formal setting permits
  reasoning over the compatibility set $\mathcal{K}_{compat}(D;E,\Pi)$.
  \item \textbf{Construction class.}
  Two distinct internal structures (two distinct continua
  $K^{(1)}_{EA} \neq K^{(2)}_{EA}$ admissible in the formal setting) such that
  the \emph{same} observation trace is admissible for both:
  \[
    D \subseteq \Sigma_{obs}(K^{(1)}_{EA},E)
    \quad\text{and}\quad
    D \subseteq \Sigma_{obs}(K^{(2)}_{EA},E).
  \]
  Equivalently, a witness that
  $\lvert \mathcal{K}_{compat}(D;E,\Pi) \rvert > 1$ holds for the claimed
  diagnosability regime.
  \item \textbf{Falsified proposition (if such constructions exist).}
  Any proposition asserting \emph{uniqueness} of structural inference from $D$
  under the stated $(E,\Pi)$ (i.e., $\lvert\mathcal{K}_{compat}(D;E,\Pi)\rvert=1$
  or an equivalent uniqueness claim) is refuted by a single CE-1 witness.
\end{enumerate}

\paragraph{CE-2 (STOP false positive under admitted states).}
\begin{enumerate}
  \item \textbf{Assumptions.}
  A STOP predicate is used as an epistemic classifier (not a recommendation) and
  is intended to hold only when diagnosis is not defined or degenerates under
  the paper's formal semantics and admissibility assumptions.
  \item \textbf{Construction class.}
  An admissible situation in which STOP is asserted, yet there exists at least
  one compatible realization supporting non-degenerate epistemic diagnosis within
  the same formal setting:
  \[
    \text{STOP}(D;E,\Pi) = \text{true}
    \quad\text{and}\quad
    \exists K \in \mathcal{K}_{compat}(D;E,\Pi)\ \text{such that the claimed
    diagnostic statement is well-defined for } K.
  \]
  (No algorithm is required; the construction is existential.)
  \item \textbf{Falsified proposition (if such constructions exist).}
  Any proposition asserting \emph{soundness} of the STOP predicate in the sense
  “STOP implies diagnosability loss under the formal semantics” is refuted.
\end{enumerate}

\paragraph{CE-3 (STOP false negative under diagnosability loss).}
\begin{enumerate}
  \item \textbf{Assumptions.}
  The paper's formal scope includes an admissible state space and a notion of
  diagnosability loss (e.g., feasibility collapse, non-identifiability to the
  point that no invariant diagnostic statement remains).
  \item \textbf{Construction class.}
  A situation in which STOP is \emph{not} asserted, yet diagnosability has
  degenerated in the relevant sense under the same formal assumptions:
  \[
    \text{STOP}(D;E,\Pi) = \text{false}
    \quad\text{and}\quad
    \forall \text{ candidate diagnostic statements } Q \text{ in scope, }
    Q \text{ is not invariant over } \mathcal{K}_{compat}(D;E,\Pi),
  \]
  or equivalently, the formal prerequisites for meaningful diagnosis fail while
  STOP remains false.
  \item \textbf{Falsified proposition (if such constructions exist).}
  Any proposition asserting \emph{completeness} of the STOP predicate with
  respect to the stated notion of diagnosability loss is refuted.
\end{enumerate}

\paragraph{CE-4 (Phase assignment non-invariance under admissible equivalences).}
\begin{enumerate}
  \item \textbf{Assumptions.}
  Phase logic assigns phase labels/regions as epistemic classifications that are
  intended to be representation-invariant under admissible equivalences
  (e.g., relabelings or semantics-preserving encodings within the formal
  setting).
  \item \textbf{Construction class.}
  Two representations of the \emph{same} admitted underlying structure (equivalent
  in the intended sense of the formal setting) that yield different phase
  assignments:
  \[
    K^{(1)}_{EA} \equiv K^{(2)}_{EA}
    \quad\text{but}\quad
    \text{Phase}(K^{(1)}_{EA};E,\Pi) \neq \text{Phase}(K^{(2)}_{EA};E,\Pi).
  \]
  \item \textbf{Falsified proposition (if such constructions exist).}
  Any proposition claiming that phase labels are well-defined on equivalence
  classes (i.e., representation-invariant) is refuted by a single CE-4 witness.
\end{enumerate}

\paragraph{Remark (status of falsification).}
These templates define \emph{logical refuters} for proposition types used in the
paper. They do not prescribe how to find such constructions, and they do not
claim empirical prevalence. The falsification mode is epistemic: existence of a
single admissible counterexample suffices.
